\setcounter{section}{42}
\section{Propiedad Intelectual}

Nombre completo: La protección jurídica de los programas de ordenador. Los medios de comprobación de la legalidad y control del software.

Apuntes de clase.

\localtableofcontents

\subsection{Historia}
    Tradicionalmente se ha vendido la máquina y se regalaba el software. Con los aos se ha cambiado y se ha protegido la propiedad intelectual. Con el RD 1/1996`se refundan las normas relativas a propiedad intelectual. En sucesivos cambios se introdujo el Canon Digital y se eliminaría por una tasa compensatoria dentro de los PGE.

    La ley no define, pero si impone qué requisitos: ha de ser original y estar plasmada en un medio material. La ley incluye tres tipos de obras: las creaciones originales, que incluyen los programas de ordenador, las obras derivadas de otros sin prejuicio de su autor original como traducciones o adaptaciones, y las obras con derechos afines conexos. La ley expresamente excluye textos como leyes, sentencias, deliberaciones de organismos públicos ni las traducciones, de poder generar derechos.

    La autoría individual se otorga a la persona natural que crea una obra. De esta se puede beneficiar una empresa. Cuando confluyen varias personas puede haber problemas determinando su autoría. Se suele tomar la forma de autoría unitaria o colaborativa (cada autor lo es de una parte) o colectiva. En la colaborativa, cada autor lo es de una parte distinta que se puede separar, explotar y cobrar individualmente, mientras que en las colectivas no se puede separar y las decisiones se deben decidir colectivamente. La ley considera al autor a la persona que edita y divulga (natural o jurídica) salvo pacto en contrario.

    El autor posee los derechos de una obra, pueden ser morales y patrimoniales. El derecho moral es irrenunciable, inalienable y perpétuo, incluye la divulgación, el formato y el reconocimiento de la condición de autor de la obra, y de la integridad. También tiene derecho de retirarla o modificarla. Los derechos patrimoniales si son renunciables y heredables, le corresponde al autor en exclusiva, el cual puede vender. Incluye los derechos de reproducción, distribución, comunicación pública y transformación.

    La producción de copias se limita a la copia privada, que se permite y compensa. Los materiales capaces de hacer copia privada están gravados, salvo excepciones.

    Los derechos patrimoniales no son eternos. Si el autor es una persona natural son 70 aos desde la muerte (por aos completos). En las obras colectivas son 70 años desde el siguiente 1 de enero de su divulgación.

    El título VII recoge los programas informáticos. La definición es vaga, pero incluye la documentación preparatoria, pero no la técnica. Tampoco protege las ideas o principios en los que se basan. Los programas no están protegidos por patentes. Si se desarrolla en conjunto HW y SW si que está protegido por patentes. Los virus y programas maliciosos no están protegidos. Existe la excepción en los autores que crean bajo contrato, en cuyo caso la autoría cae en la empresa.

    La administración pide todos los derechos que contrata, menos los morales. En los derechos de explotación, la venta del programa agota el derecho, por eso no se vende, se licencia. 

    Se consideran cuatro límites al derecho de explotación que son los de stallman.

    En España hay un único registro de la propiedad intelectual, delegado a las comunidades, el registro no es obligatorio. Se einscriben derechos propios y cedidos, son declarativas, pero no garantistas. Se presumen correctos salvo prueba en contratio.